\chapter{Conclusions and Future Work}

When I started planning this work, I was clear that I wanted to build something that actually worked, not just a well-explained set of concepts on paper. And I believe I have achieved that: I have a running Kubernetes cluster with security policies that block dangerous configurations before they ever execute, that automatically correct what they can correct, integrated with ArgoCD so that deployments are declarative and traceable, and with a pipeline that analyses images for vulnerabilities before publishing anything. All connected, all working. That gives me some satisfaction.

What this work has given me most is not, in truth, any of the specific tools. It has been understanding how DevSecOps teams really work: that security is not a layer you add at the end, but something you have to keep in mind from the very beginning. You see that a lot in the theory, but until you set it up yourself, with all the mistakes included, you do not really internalise it. I learned by doing, and many things I only understood because I ran into them.

The finding that surprised me most, and that I am most pleased with, came when setting up the vulnerability analysis pipeline. I had two tools working on the same image: Trivy and Grype. Trivy passed it. Grype found an important vulnerability in the Python image and blocked the pipeline. Two tools, one image, different results. I could have lowered Grype's threshold so the image would pass and put the problem aside, but I did not: I went to the source, updated the base image to a more recent version of Python, and the vulnerability disappeared. Seeing that in practice taught me more than any article about defence in depth. It makes sense to use multiple layers of security, not because a book says so, but because I saw it happen.

I also found it interesting to discover how ArgoCD and Kyverno complement each other to keep the system in the correct state. I watched Kyverno correct insecure configurations before they reached validation, and ArgoCD detect and revert any manual change to the cluster. There is an elegant logic in how these two tools fit together. I also learned that if Kyverno fails or runs out of resources, it blocks the entire cluster---because it is designed to fail safely, not to let requests through without review. I learned that the hard way: the node's disk filled up, Kyverno crashed, and the cluster was blocked. It was a rough moment. But it also taught me something that does not appear in the manuals: that these tools have real consequences when something goes wrong, and that resource management is part of the job.

I am aware that some things were left out. I did not manage to integrate Sysdig Secure to monitor what happens in the cluster while applications are running. The comparison with OPA Gatekeeper I did at the code level, writing the equivalent policies in Rego, but without deploying Gatekeeper in the cluster, because having two admission controllers active at the same time would have been messy. And I did not implement image signing with Cosign either, which would have added another interesting security layer to the chain. These are things I know are missing, and that I would address another time, but I also believe the work has a reasonable scope for what it is.

Looking ahead, I have several things I would like to explore. One is genuinely integrating a runtime security tool in a cloud environment, which is where it makes sense. Another is automating the entire environment setup with Ansible, to be able to reproduce it without doing anything by hand, which is how any serious infrastructure should work. And I would also like to set up the same pipeline in GitLab CI instead of GitHub Actions, because GitLab is widely used in companies and I am curious to see the differences.

This work is not an end point. It is more like the beginning of something. I know there is still a lot to learn, that there are parts I could have done better, and that the world of Kubernetes security moves fast. But I come out of it with a foundation I did not have before, with the feeling of having truly understood how these tools work, and with the motivation to keep going. And for me, that is enough.
